\documentclass[11pt,a4paper]{report}

\usepackage[utf8]{inputenc}
\usepackage[T1]{fontenc}
\usepackage{lmodern}
\usepackage{microtype}
\usepackage{geometry}
\geometry{margin=1in, headheight=14pt}
\usepackage{amsmath,amssymb,amsfonts,bm}
\usepackage{booktabs,longtable,array,tabularx}
\usepackage{graphicx}
\usepackage{xcolor}
\usepackage{hyperref}
\usepackage{fancyhdr}
\usepackage{titlesec}
\usepackage{enumitem}
\usepackage{listings}
\usepackage[most]{tcolorbox}

% Color definitions
\definecolor{primarynavy}{RGB}{11, 37, 69}
\definecolor{secondaryblue}{RGB}{19, 64, 116}
\definecolor{accentteal}{RGB}{17, 94, 89}
\definecolor{codegray}{RGB}{245, 247, 250}
\definecolor{codeframe}{RGB}{203, 213, 225}
\definecolor{textdark}{RGB}{15, 23, 42}
\definecolor{alertred}{RGB}{185, 28, 28}
\definecolor{successgreen}{RGB}{21, 128, 61}

% Hyperref configuration
\hypersetup{
    colorlinks=true,
    linkcolor=secondaryblue,
    citecolor=accentteal,
    urlcolor=secondaryblue,
    pdfauthor={scienceStack Engineering Group},
    pdftitle={scienceStack: Rigorous Scientific Engineering Workflows for Google Antigravity},
    pdfsubject={Scientific Computing, Autonomous AI Agents, Equation Verification},
    pdfkeywords={Google Antigravity, scienceStack, numerical stability, equation provenance, golden benchmarks}
}

% Code Listing configuration
\lstset{
    backgroundcolor=\color{codegray},
    basicstyle=\small\ttfamily\color{textdark},
    breaklines=true,
    breakatwhitespace=true,
    frame=single,
    rulecolor=\color{codeframe},
    numbers=left,
    numberstyle=\tiny\color{gray},
    numbersep=8pt,
    captionpos=b,
    keywordstyle=\color{secondaryblue}\bfseries,
    commentstyle=\color{accentteal}\itshape,
    stringstyle=\color{alertred},
    showstringspaces=false,
    tabsize=2
}

% Page styling
\pagestyle{fancy}
\fancyhf{}
\fancyhead[L]{\small\color{secondaryblue}\textsf{\nouppercase{\leftmark}}}
\fancyhead[R]{\small\color{secondaryblue}\textsf{scienceStack Manual (v0.22.0)}}
\fancyfoot[C]{\thepage}
\renewcommand{\headrulewidth}{0.5pt}

% Section title formats
\titleformat{\chapter}[display]
  {\normalfont\huge\bfseries\color{primarynavy}}{\chaptertitlename\ \thechapter}{16pt}{\Huge}
\titleformat{\section}
  {\normalfont\Large\bfseries\color{primarynavy}}{\thesection}{1em}{}
\titleformat{\subsection}
  {\normalfont\large\bfseries\color{secondaryblue}}{\thesubsection}{1em}{}
\titleformat{\subsubsection}
  {\normalfont\normalsize\bfseries\color{accentteal}}{\thesubsubsection}{1em}{}

% Custom tcolorbox environments
\newtcolorbox{principlebox}[1]{
    colback=accentteal!5!white,
    colframe=accentteal,
    fonttitle=\bfseries\sffamily,
    title={Scientific Principle: #1},
    enhanced,
    attach boxed title to top left={yshift=-2mm, xshift=4mm},
    boxed title style={colback=accentteal},
    sharp corners,
    boxrule=1pt,
    top=8pt,
    bottom=8pt
}

\newtcolorbox{directivebox}[1]{
    colback=primarynavy!5!white,
    colframe=primarynavy,
    fonttitle=\bfseries\sffamily,
    title={Antigravity Directive: #1},
    enhanced,
    attach boxed title to top left={yshift=-2mm, xshift=4mm},
    boxed title style={colback=primarynavy},
    sharp corners,
    boxrule=1pt,
    top=8pt,
    bottom=8pt
}

\newtcolorbox{warningbox}[1]{
    colback=alertred!5!white,
    colframe=alertred,
    fonttitle=\bfseries\sffamily,
    title={Numerical Safety Alert: #1},
    enhanced,
    attach boxed title to top left={yshift=-2mm, xshift=4mm},
    boxed title style={colback=alertred},
    sharp corners,
    boxrule=1pt,
    top=8pt,
    bottom=8pt
}

\newtcolorbox{notebox}[1]{
    colback=secondaryblue!5!white,
    colframe=secondaryblue,
    fonttitle=\bfseries\sffamily,
    title={Architectural Note: #1},
    enhanced,
    attach boxed title to top left={yshift=-2mm, xshift=4mm},
    boxed title style={colback=secondaryblue},
    sharp corners,
    boxrule=1pt,
    top=8pt,
    bottom=8pt
}

\begin{document}

\begin{titlepage}
    \centering
    \vspace*{1.5cm}
    {\Huge \bfseries \color{primarynavy} scienceStack}\\[0.4cm]
    {\Large \bfseries \color{secondaryblue} Rigorous Scientific Engineering Workflows}\\[0.2cm]
    {\large \textsf{Operational Manual for the Google Antigravity Plugin}}\\[1.5cm]
    \rule{\linewidth}{1.5pt}\\[0.6cm]
    {\large \textbf{Plugin Specification \& Operational Guide}}\\[0.2cm]
    {\normalsize Version 0.22.0 --- October 2026}\\[0.6cm]
    \rule{\linewidth}{1.5pt}\\[2.5cm]
    \begin{minipage}{0.48\textwidth}
        \begin{flushleft} \large
            \textbf{Components:}\\
            12 Skills\\
            6 Playbooks, 6 Principles\\
            2 Subagents\\
            4 Lifecycle Hooks\\
            1 Directive Rule File
        \end{flushleft}
    \end{minipage}
    \hfill
    \begin{minipage}{0.48\textwidth}
        \begin{flushright} \large
            \textbf{Domains:}\\
            PDEs, CFD, orbital mechanics\\
            Stiff ODE kinetics\\
            Numerical linear algebra
        \end{flushright}
    \end{minipage}
    \vfill
    {\small \textsf{Adapted from the \texttt{pstack} architecture by Lauren Tan (poteto) at Cursor.}\\
    \textsf{Specialized for equation provenance, dimensional analysis, and numerical verification in Google Antigravity.}}\\[0.8cm]
\end{titlepage}

\pagenumbering{roman}
\tableofcontents
\newpage

\section*{Executive Summary}
\addcontentsline{toc}{chapter}{Executive Summary}

Generated code that satisfies local tests can still be wrong in the ways that matter to computational science: a flipped sign in a diffusion operator, an off-by-one in a stencil, or an explicit matrix inverse that squares the condition number. \textbf{scienceStack} packages a small set of enforceable rules, workflows, and tools for Google Antigravity that target exactly these failures. It is organized around six principles:
\begin{enumerate}[itemsep=2pt]
    \item \textbf{Equation provenance}: every implemented formula carries its continuous LaTeX form, a literature citation, units, and sign conventions.
    \item \textbf{Dimensional discipline}: both sides of every equation have the same dimensions; arguments of transcendental functions are dimensionless.
    \item \textbf{Numerical stability first}: no explicit inversion, guarded denominators, and \texttt{log1p}/\texttt{expm1} for small arguments.
    \item \textbf{Conservation laws}: Hamiltonian systems use symplectic integrators; energy, momentum, and angular momentum are tracked in tests.
    \item \textbf{Physics-modular architecture}: pure physical laws live in stateless modules, decoupled from time steppers, meshes, and I/O.
    \item \textbf{Analytical limits and manufactured solutions}: solvers are verified against closed-form solutions or the Method of Manufactured Solutions with explicit error norms.
\end{enumerate}

\newpage
\pagenumbering{arabic}

\chapter{Installation and How to Use}

\section{Installing the Plugin}
The plugin complies with Agent Plugins 1.0.0; its root manifest is \texttt{plugin.json}.
\begin{lstlisting}[language=bash]
# 1. Validate compliance with Agent Plugins 1.0.0 (Strict mode)
agy plugin validate /home/seo3919/AI_skills/scienceStack_antigravity --strict

# 2. Install globally in Antigravity
agy plugin install /home/seo3919/AI_skills/scienceStack_antigravity

# Or workspace-local:
# Place scienceStack_antigravity in .agents/plugins/sciencestack
\end{lstlisting}
Verify installation status:
\begin{lstlisting}[language=bash]
agy plugin list
# - sciencestack (enabled, 0.22.0)
#   Components: 12 skills, 2 agents, 1 rule, 4 hooks
\end{lstlisting}

\section{Skills}
Skills can be triggered via interactive slash commands or dynamically by Antigravity agents:
\begin{itemize}[itemsep=2pt]
    \item \texttt{/science-mode <task>}: master workflow. Selects exactly one of six playbooks and enforces the six principles. User-invoked.
    \item \texttt{/setup-sciencestack}: writes per-role model choices to configuration. User-invoked.
    \item \texttt{astronomy-computing}: Kepler solvers, symplectic N-body dynamics, SO(3) coordinate frames.
    \item \texttt{accelerated-scientific-computing}: memory layout, SIMD vectorization, and precision guards.
    \item \texttt{check-against-reference}: runs the user's solver on Kepler's equation, Lane-Emden, the Sod shock tube, or Robertson kinetics, or compares its output with golden data, through \texttt{scripts/check\_against\_reference.py}. The script decides (exit status 0 pass, 1 fail, 2 nothing checked, 3 verdict withheld or manifest weakened); \texttt{convergence-ode} and \texttt{convergence-bvp} measure observed order of convergence; \texttt{--repeat} detects nondeterminism; every run appends to a SHA-256 bound evidence log (\texttt{.verification/log.jsonl}); tolerance loosening after failure on unchanged files is withheld (exit 3) unless justified; \texttt{run-manifest} re-runs committed manifests.
    \item \texttt{run-tests}: runs unittest or pytest through \texttt{scripts/run\_tests.py}, extracts exact counts from framework internals, fails when nothing ran or everything was skipped, and records results bound to the worktree hash; \texttt{run\_tests.py fresh} confirms whether tests passed on current files.
    \item \texttt{test-ratchet}: \texttt{scripts/test\_ratchet.py --baseline-ref REF} parses ASTs of test files against a git baseline and exits 3 on unapproved looser tolerance, removed tests, fewer assertions, new skip/xfail, or fewer parametrize cases; approvals go in \texttt{test-ratchet-approvals.json}.
    \item \texttt{numerical-stability-audit}: audits explicit matrix inversion, cancellation, and unseeded random number generator states via \texttt{scripts/audit\_numerical\_ast.py}.
    \item \texttt{precision-audit}, \texttt{scientific-equation-verifier}, \texttt{golden-dataset-benchmark}, \texttt{benchmark-checklist}: audit and verification skills loaded on demand.
\end{itemize}

\section{Runtime Directives}
Directives are defined in \texttt{rules/AGENTS.md}: the core directives (provenance, stability, symplectic integration, modularity, reference checks, test ratchet, evidence logs) and absolute script locations.

\chapter{Architecture}

\section{Subagents and Models}
\begin{table}[h]
\centering
\small
\begin{tabular}{llll}
\toprule
\textbf{Agent} & \textbf{Model} & \textbf{Tools} & \textbf{Role} \\
\midrule
\texttt{science-master} & \texttt{flash} / \texttt{inherit} & all & kernels, MMS tests, refactors, benchmark runners \\
\texttt{theory-reviewer} & \texttt{pro} & read-only & line-by-line audit of equations, signs, index order, invariants \\
\bottomrule
\end{tabular}
\caption{Subagents. \texttt{theory-reviewer} operates in a read-only context.}
\end{table}

\section{Where Things Live}
\begin{lstlisting}[language=bash]
plugin.json                     Agent Plugins 1.0.0 manifest
hooks/hooks.json                PreToolUse, PostToolUse, and StopGate hooks
rules/AGENTS.md                 Antigravity runtime directives
agents/                         science-master, theory-reviewer
references/                     compact delegate principles (read by science-master)
skills/science-mode/            playbooks/ (6) and principles/ (6)
skills/                         check-against-reference, run-tests, test-ratchet, ...
scripts/                        hooks, AST auditor, MMS and Jacobian generators
benchmarks/                     analytic reference solvers
tests/                          unit and integration tests
\end{lstlisting}

\chapter{The Six Principles}

\begin{principlebox}{Equation Provenance}
Document the governing equation in the docstring, for example
\[ \frac{\partial u}{\partial t} = \alpha \nabla^2 u \quad [\mathrm{SI}], \]
with its literature source, coordinate system, and sign convention. The AST auditor reports scientific functions whose docstring has no LaTeX equation.
\end{principlebox}

\begin{principlebox}{Dimensional Discipline}
For every equation $[\mathrm{LHS}] = [\mathrm{RHS}]$. Arguments of $\exp$, $\ln$, $\sin$ must be dimensionless. Non-dimensionalize coupled systems with characteristic scales.
\end{principlebox}

\begin{principlebox}{Numerical Stability First}
For a linear system $A\mathbf{x}=\mathbf{b}$ the relative error obeys
\[ \frac{\|\delta \mathbf{x}\|}{\|\mathbf{x}\|} \;\lesssim\; \kappa(A)\left(\frac{\|\delta A\|}{\|A\|}+\frac{\|\delta \mathbf{b}\|}{\|\mathbf{b}\|}\right). \]
Use \texttt{solve(A, b)} or a factorization rather than forming $A^{-1}$. Replace $\ln(1+x)$ by \texttt{log1p}, $e^x-1$ by \texttt{expm1}, $1-\cos x$ by $2\sin^2(x/2)$.
\end{principlebox}

\begin{principlebox}{Conservation Laws}
Hamiltonian systems $H(\mathbf{q},\mathbf{p})$ use symplectic integrators (kick-drift-kick leapfrog, Yoshida splitting), which conserve a nearby shadow Hamiltonian $\widetilde H = H + O(\tau^2)$ and show no secular energy drift. Track $E(t)$, $\mathbf{P}(t)$, $\mathbf{L}(t)$ in tests with an explicit tolerance. Explicit Runge--Kutta is not used for long-term orbit integration.
\end{principlebox}

\begin{principlebox}{Physics-Modular Architecture}
Forces, equations of state, and rate equations are pure, stateless functions in a \texttt{physics/} module that other solvers can import without a time stepper, mesh, or I/O.
\end{principlebox}

\begin{principlebox}{Analytical Limits and MMS}
For a differential operator $\mathcal{L}$, choose a smooth $u_{\mathrm{mms}}$, compute $f=\mathcal{L}[u_{\mathrm{mms}}]$ symbolically, solve $\mathcal{L}[u]=f$ numerically and check that the error norm falls with the expected order $p$ under grid refinement:
\[ p_{\mathrm{obs}} = \frac{\ln(e_h/e_{h/2})}{\ln 2}. \]
\end{principlebox}

\chapter{Playbooks}
\texttt{/science-mode} selects and executes the specific playbook matching the task.
\begin{table}[h]
\centering
\small
\begin{tabular}{lp{10cm}}
\toprule
\textbf{Playbook} & \textbf{Purpose} \\
\midrule
\texttt{equation-implementation} & Turn a published equation into verified code with LaTeX docstrings, dimension checks, and limit tests. \\
\texttt{numerical-verification} & Truncation-error order, grid-refinement convergence, conservation stability. \\
\texttt{perf-issue} & Profile and remove a bottleneck while asserting $L_2$ parity with the previous output. \\
\texttt{hillclimb} & Iterate against a quantitative target: one change, one measurement, keep or revert, with a decision log. \\
\texttt{bug-fix} & Reproduce on a minimal analytic case, find the root cause by evidence, failing test first. \\
\texttt{refactoring} & Pin behavior with a characterization test, then restructure into stateless physics modules. \\
\bottomrule
\end{tabular}
\end{table}

\chapter{Safety Hooks and Tools}

\section{PreToolUse Safety Gate}
\texttt{scripts/hook\_pre\_tool\_guard.py} intercepts tool writes (\texttt{write\_file}, \texttt{edit\_file}, and terminal write invocations) and blocks proposed code that contains three critical hazards:
\begin{enumerate}[itemsep=2pt]
    \item \textbf{Explicit matrix inversion}: \texttt{numpy.linalg.inv}, \texttt{scipy.linalg.inv}, \texttt{torch.linalg.inv}, \texttt{torch.inverse}, and JAX/CuPy/SciPy-sparse equivalents. Names are resolved through AST import aliases.
    \item \textbf{Truncated $\pi$}: literal values \texttt{3.14159} and \texttt{3.1416}.
    \item \textbf{Mutable NumPy default argument}: \texttt{def f(x=np.zeros(3))}.
\end{enumerate}
Comments, docstrings, and strings are parsed away by the AST and never trigger the gate.

\section{Building Both Sides of a Comparison}
A manifest \texttt{compare} check can produce its reference and/or its test output by running programs instead of trusting committed files: \texttt{ref\_from} and \texttt{test\_from} give argument lists (never a raw shell), the output (\texttt{stdout} or file path), and \texttt{pinned}, the SHA-256 of every source and input file. A changed pinned file fails the check before anything runs. For \texttt{ref\_from}, the committed reference must reproduce the rebuilt numbers within \texttt{reproduce\_rtol} and \texttt{reproduce\_atol}.

\section{Verification Hooks}
In a git project that uses scienceStack verification (contains \texttt{reference-checks.json}, \texttt{test-ratchet-approvals.json}, or \texttt{.verification/}):
\begin{itemize}[itemsep=2pt]
    \item \texttt{scripts/hook\_test\_guard.py} (PreToolUse): Compares proposed edits against baseline ASTs and denies/asks if a test is weakened (looser tolerance, removed test/assertion, new skip/xfail, fewer parametrizations).
    \item \texttt{scripts/hook\_record\_edit.py} (PostToolUse): Records files modified during the active session.
    \item \texttt{scripts/hook\_stop\_gate.py} (StopGate): Blocks completion once if edited solvers or compared outputs lack fresh passing checks, recorded tests are stale, or edited tests lack fresh ratchet runs.
\end{itemize}
Hooks fail closed: internal crashes or unexpected formats deny modifications until addressed. Set \texttt{SCIENCESTACK\_HOOKS=off} in the environment to temporarily bypass verification hooks.

\begin{warningbox}{Limits}
The gate cannot see an inversion built dynamically (\texttt{getattr}), \texttt{np.matrix} objects used through \texttt{.I}, or code written by unhooked external processes. It is a safety net, not an absolute formal proof.
\end{warningbox}

\section{AST Auditor}
\texttt{python3 scripts/audit\_numerical\_ast.py <file|dir> [--strict] [--json]} reports:
\texttt{BAN\_MATRIX\_INV} (error); \texttt{CANCELLATION\_LOG1P}, \texttt{CANCELLATION\_EXPM1}, \texttt{CANCELLATION\_ONE\_MINUS\_COS}, \texttt{PRECISION\_FLOAT32}, \texttt{TRUNCATED\_PI}, \texttt{MUTABLE\_ARRAY\_DEFAULT} (warnings); \texttt{MISSING\_LATEX\_PROVENANCE} (info).

\section{MMS and Jacobian Generators}
\begin{lstlisting}[language=bash]
python3 scripts/mms_generator.py --exact "sin(pi*x)*exp(-t)" \
    --operator "diff(u, t) - alpha*diff(u, x, 2)" --params "alpha=0.1" --output fixture.py
python3 scripts/derive_jacobian.py --vars "x,y" --eqs "x**2 + y, sin(x)*y"
\end{lstlisting}
For $u=e^{-t}\sin\pi x$ and $\mathcal L[u]=u_t-\alpha u_{xx}$ the generator derives the forcing
\[ f(x,t) = (\alpha\pi^2-1)\,e^{-t}\sin\pi x, \]
and writes a fixture whose convergence test is ready for solver attachment. The Jacobian generator emits \texttt{compute\_jacobian(\_state)} and analytic symbolic entries.

\section{Reference Benchmarks}
\begin{table}[h]
\centering
\small
\begin{tabular}{p{4.6cm}p{6cm}p{3.6cm}}
\toprule
\textbf{Script} & \textbf{Reference} & \textbf{Source} \\
\midrule
\texttt{cfd/exact\_riemann.py} & Sod star state $p_*=0.30313$, $u_*=0.92745$ ($\gamma=1.4$) & Toro (2009), Table 4.1 \\
\texttt{astronomy/lane\_emden\_solver.py} & $\xi_1=\pi$ for $n=1$; $3.65375$ for $n=1.5$; $6.89685$ for $n=3$ & standard polytrope values \\
\texttt{astronomy/kepler\_solver.py} & conserved $E$ and $\mathbf L$ & analytic invariants \\
\texttt{astronomy/symplectic\_nbody.py} & bounded energy error, conserved $\mathbf P$, $\mathbf L$ & symplectic structure \\
\texttt{astronomy/coordinate\_transforms.py} & $R^TR=I$, $\det R=1$ & SO(3) \\
\texttt{kinetics/robertson\_reference.py} & $y_1+y_2+y_3=1$ & mass conservation \\
\bottomrule
\end{tabular}
\caption{Benchmarks and what each is checked against.}
\end{table}

\chapter{Testing}
\begin{lstlisting}[language=bash]
pytest tests/
\end{lstlisting}
The 594-test suite verifies safety hooks (import aliases, formatting resilience), the AST auditor, MMS and Jacobian generators, reference benchmarks, convergence order calculations, ratchet mechanics, and hash-chained verification evidence ledgers.

\chapter*{Recommended Literature \& References}
\addcontentsline{toc}{chapter}{Recommended Literature \& References}
\begin{enumerate}[itemsep=6pt]
    \item Trefethen, L. N., \& Bau III, D. (1997). \textit{Numerical Linear Algebra}. SIAM.
    \item Golub, G. H., \& Van Loan, C. F. (2013). \textit{Matrix Computations}, 4th ed. Johns Hopkins University Press.
    \item Higham, N. J. (2002). \textit{Accuracy and Stability of Numerical Algorithms}, 2nd ed. SIAM.
    \item Hairer, E., Lubich, C., \& Wanner, G. (2006). \textit{Geometric Numerical Integration}, 2nd ed. Springer.
    \item Roache, P. J. (2009). \textit{Fundamentals of Verification and Validation}. Hermosa Publishers.
    \item Salari, K., \& Knupp, P. (2000). \textit{Code Verification by the Method of Manufactured Solutions}. Sandia National Laboratories, SAND2000-1444.
    \item Toro, E. F. (2009). \textit{Riemann Solvers and Numerical Methods for Fluid Dynamics}, 3rd ed. Springer.
    \item Hairer, E., \& Wanner, G. (1996). \textit{Solving Ordinary Differential Equations II}, 2nd ed. Springer.
\end{enumerate}

\end{document}
